\section{The classical level-three family in every even degree}
\label{sec:gl2-uniform}

Put $\chi=\chi_{-3}$ and let $e\geq2$ be even. The complete
classical space is $S_e=S_{2e+1}(\Gamma_0(3),\chi)$.
Its dimension and integral basis are already supplied by
K17-GL2-GENERAL, \texttt{gl2\_general\_section.tex}, (G1)--(G5).
They are cited here, with the empty-space convention at $e=2$:
\[
 r_e=\dim S_e=\left\lfloor\frac{2e-2}{3}\right\rfloor,\qquad
 f_j=D a^{\,2e-5-3j}b^j\quad(0\leq j<r_e),
 \tag{GL.1}
\]
where
\[
 a=\sum_{\alpha\in\mathbf Z[\omega]}q^{N\alpha},\qquad
 b=\frac{\eta(3\tau)^9}{\eta(\tau)^3},\qquad
 D=\eta(\tau)^6\eta(3\tau)^6 .
\]
The primary ring statement is Cl\'ery--van der Geer,
\emph{Generating Picard modular forms by means of invariant theory},
arXiv:2110.00849, paragraph \emph{Modular forms on $\Gamma_1(3)$},
printed p.~12; the same work, Remark~15.2, printed p.~33,
gives the weight-six cusp form. The basis begins
$f_j=q^{j+1}+\cdots$ and is integral and unitriangular.

\begin{proposition}[KT-GL2-CM]
The CM subspace of $S_e$ has dimension
\[
 c_e=\begin{cases}1,&3\mid e,\\0,&3\nmid e.\end{cases}
 \qquad
 \dim S_e^{\mathrm{nonCM}}=r_e-c_e.
 \tag{GL.2}
\]
When $3\mid e$, its unique normalized eigenform is the theta series
of the conductor-one Hecke character
$\psi_e((\alpha))=\alpha^{2e}$ of $\mathbf Q(\omega)$.
At a split prime $p=\pi\bar\pi$ its eigenvalue is
\[
             a_p(\psi_e)=\pi^{2e}+\bar\pi^{2e}.
 \tag{GL.3}
\]
At an inert prime its eigenvalue is zero.
\end{proposition}

\begin{proof}
The nebentypus has primitive conductor three. Therefore the space
has no oldforms from a proper divisor of the level, and every
normalized eigenform has primitive level three.
The CM classification says that such a form comes from an
algebraic Hecke character of an imaginary quadratic field, with
level the absolute field discriminant times the norm of its
conductor. We use exactly Sch\"utt,
\emph{CM newforms with rational coefficients},
arXiv:math/0511228v5, Theorem~1.1 (Ribet's classification),
printed p.~2, and Theorem~1.4 (Hecke--Shimura),
printed p.~3. The cited classification is not restricted to
rational coefficient fields.

The level formula forces the field to be $\mathbf Q(\omega)$
and the conductor to be the unit ideal. Its class number is one.
A character with infinity type $2e$ and conductor one exists
precisely when $\alpha^{2e}$ is unchanged on replacing a generator
of a principal ideal by a unit multiple. Since its units are
$\mu_6$, this means $6\mid2e$, equivalently $3\mid e$.
When this holds, class number one makes the character unique.
The character with conjugate infinity type gives the same theta
series, by reindexing ideals by conjugation, so it gives no second
eigenform. The two prime ideals over a split prime give (GL.3).
There is no ideal of norm an inert rational prime, giving zero
there. Hecke--Shimura supplies a cuspidal eigenform in exactly
the stated level and nebentypus. This proves both existence and
exhaustion. Subtracting its dimension from (GL.1) proves (GL.2).
\end{proof}

\subsection{A uniform exact determinant formula}

For $p\ne3$ let
\[
 L_{ij}=[q^i]f_j,\quad
 B_{ij}=[q^{pi}]f_j+
       {\bf1}_{p\mid i}\chi(p)p^{2e}[q^{i/p}]f_j
 \quad(1\leq i\leq r_e,\ 0\leq j<r_e).
\]
Then, by the already proved integral basis and the classical
Hecke formula,
\[
       Q_{e,p}(X)=\det(XI-L^{-1}B)\in\mathbf Z[X].
 \tag{GL.4}
\]
This gives every eigenvalue, with multiplicity, uniformly in $e$.
It is the matrix form of K17-GL2-GENERAL, not a new modularity
theorem. The required Fourier coefficients extend only through
degree $pr_e$.

The fresh computation \texttt{evidence/gl2\_uniform.py}
constructs these matrices for every even $2\leq e\leq30$
and $p=7,13$. It expands the eta products by their logarithmic
derivatives. Explicitly, if
$F(q)=\prod_{n\geq1}(1-q^n)^s(1-q^{3n})^t
=\sum_{n\geq0}c_nq^n$, then
\[
 c_0=1,\qquad
 n c_n=\sum_{j=1}^n
   \bigl(-s\sigma_1(j)-3t\,{\bf1}_{3\mid j}\sigma_1(j/3)\bigr)c_{n-j}.
 \tag{GL.5}
\]
This follows by multiplying $qF'/F$ by $F$ and comparing
coefficients. Every division by $n$ is checked to be exact.
This implementation uses (GL.5), rather than the predecessor's
binomial expansion of the eta products.
It checks further Hecke coefficients after the identifying
ones and compares the characteristic-polynomial evaluation
with a direct determinant. The complete higher4 matrices at
$e=26,28,30$ agree entry by entry with the attached K17 data.
The known $e=6$ CM eigenvalue is independently recovered as
$-153502$ at seven; the other two eigenvalues are $40250$.

These classical calculations settle the classical numerical half
of the three requested higher degrees. Arithmetic vanishing
also uses the full unitary comparison and source coverage.

\paragraph{Geometric interpretation.}
The eigenforms give rank-two compatible systems of pure weight
$2e$; their twists occurring in the Selmer comparison have
weight zero. The CM summand is an induced rank-one system on
$\operatorname{Spec}\mathbf Q(\omega)$, and its Euler factor
at a split prime has roots $\pi^{2e},\bar\pi^{2e}$.
The decomposition (GL.2) is a statement about these classical
automorphic systems. It does not identify the stable and
endoscopic summands of the compact unitary space.
